% TOP_PROBLEM: {"id": 3116, "title": "Star chromatic index of complete graphs", "category_id": 3, "category": {"id": 3, "name": "graph_theory", "display_name": "Graph Theory", "description": "Problems involving graphs, networks, and their properties.", "slug": "graph-theory", "order_index": 3, "created_at": "2026-07-31T15:26:25.671Z"}, "status": "open", "problem_number": "OPG-37275", "set_id": 12}
% FIRST_POSTED: 2026-10-01
% ATTEMPT_STATUS: unresolved
% UnsolvedMath ID 3116; source catalogue code OPG-37275
% !TeX program = lualatex
% GPT-6 Astra Ultra research attempt, 2026-10-01; self-contained partial work.
% Completion estimate: no numerical percentage assigned; full problem unresolved.
\documentclass[11pt,a4paper]{article}
\usepackage[margin=25mm]{geometry}
\usepackage{fontspec}
\setmainfont{lmroman10-regular.otf}[BoldFont=lmroman10-bold.otf,
 ItalicFont=lmroman10-italic.otf,BoldItalicFont=lmroman10-bolditalic.otf]
\usepackage{amsmath,amssymb,mathtools}
\usepackage{unicode-math}
\setmathfont{latinmodern-math.otf}
\AtBeginDocument{\let\setminus\smallsetminus}
\usepackage{xurl,hyperref}
\hypersetup{colorlinks=true,urlcolor=blue}
\setcounter{secnumdepth}{0}
\setlength{\parindent}{0pt}
\setlength{\parskip}{5pt}
\setlength{\emergencystretch}{3em}
\begin{document}
\section*{Star chromatic index of complete graphs}\label{3116}
{\small UnsolvedMath ID 3116 · OPG-37275 · Graph Theory · OpenGarden}
\subsection{Definitions and mathematical statement}
For $n\ge2$, let $K_n$ be the complete simple undirected graph on
$\{1,\ldots,n\}$. A proper edge colouring assigns different colours to
any two edges sharing a vertex. It is a star edge colouring if, in
addition, each path with four edges and each cycle with four edges
receives at least three colours. Paths are simple, with five distinct
vertices, and need not be induced; this distinction matters in a
complete graph.

Let $\chi'_s(K_n)$ denote the minimum size of a colour palette admitting
such a colouring. The question is whether there exists an absolute
constant $C>0$ such that
\[
 \chi'_s(K_n)\le Cn\qquad\text{for every }n\ge2.
\]
This is the precise meaning of the requested $O(n)$ upper bound.
Allowing the inequality only for all sufficiently large $n$ is equivalent,
since the finitely many smaller values can be absorbed by increasing $C$.

The palette and colouring may change with $n$, but $C$ must not grow
with $n$. Every pair of colours induces a union of paths of at most
three edges, which gives another way to express the forbidden-pattern
condition. Since all $n-1$ edges at a vertex require different colours,
a linear upper bound would also give the order of growth
$\chi'_s(K_n)=\Theta(n)$.
The problem does not ask that the same constant work for an arbitrary
graph in terms of its maximum degree; it is the family of complete
graphs that is quantified over.
\subsection{Short English statement}
Is a constant times $n$ colours enough to star edge-colour every complete graph $K_n$?
\subsection{Research attempt}
\paragraph{Scope and process.}
GPT-6 Astra Ultra, 1 October 2026. This attempt keeps the catalogue statement above, checks primary literature, and develops a restricted argument with an adversarial gap audit. The proofs below are the mathematical output of this attempt, not a claim of novelty or external certification.

\paragraph{Literature and attack plan.}
Dvořák, Mohar and Šámal [R1] give a near-linear upper bound for the star chromatic index, and a lower bound asymptotic to $2n$ for complete graphs. The present construction is weaker, but entirely explicit. It tests the additive-colouring idea, identifies exactly why ordinary integer labels fail, and repairs that defect using a progression-free set. None of the following estimates is claimed to improve [R1].

\paragraph{Additive colouring lemma.}
Let $A$ be a finite set of integers with no nonconstant three-term arithmetic progression: $a+e=2c$ with $a,c,e\in A$ implies $a=c=e$. Label the vertices of the complete graph by $A$ and colour the edge $ab$ by the integer $a+b$.

This is proper, because $a+b=a+c$ implies $b=c$. Suppose a simple path $a,b,c,d,e$ of four edges used just two colours. Properness forces alternation, so
\[
 a+b=c+d,\qquad b+c=d+e.
\]
Subtracting gives $a-c=c-e$, hence $a+e=2c$, contradicting the distinct path vertices and the progression-free property.

For a bichromatic four-cycle $a,b,c,d,a$, alternation instead gives
\[
 a+b=c+d,\qquad b+c=d+a.
\]
Subtracting yields $2(a-c)=0$. Over the integers this implies $a=c$, again impossible for a simple cycle. Thus additive colouring is a star edge colouring whenever $A$ is progression-free.

\paragraph{An explicit family and quantitative bound.}
For $k\ge1$ let
\[
 A_k=\left\{\sum_{j=0}^{k-1}\varepsilon_j3^j:
                  \varepsilon_j\in\{0,1\}\right\}.
\]
This set has $2^k$ elements. To prove it progression-free, suppose $a+e=2c$. Adding $a,e$ in base three creates no carries, since each digit sum is at most two. Doubling $c$ also creates no carries. Uniqueness of the base-three expansion gives $a_j+e_j=2c_j$ for every digit. If $c_j=0$, both other digits are zero; if $c_j=1$, both are one. Thus $a=e=c$.

Each edge sum has $k$ base-three digits in $\{0,1,2\}$, so at most $3^k$ colours occur. Choose $k=\lceil\log_2 n\rceil$ and any $n$ labels from $A_k$. Restricting a star colouring preserves its forbidden-pattern property, giving
\[
 \chi'_s(K_n)\le3^{\lceil\log_2 n\rceil}
             \le3n^{\log_2 3}. \tag{1}
\]
This is a genuine subquadratic construction, though not a linear one. Its overhead relative to $n$ grows like $n^{\log_2 3-1}$, so (1) cannot be recast as $Cn$ with absolute $C$.

\paragraph{Why consecutive labels do not work.}
Using labels $\{0,1,\ldots,n-1\}$ would need only $2n-3$ possible sums, an attractive linear palette. But for $n\ge5$, the four-edge path with vertices
\[
 0,4,1,3,2
\]
has colours $4,5,4,5$. It is a simple path in $K_n$; its many chords do not excuse it, because the problem does not restrict to induced paths. This explicitly refutes the most immediate linear additive construction.

\paragraph{Verification and remaining gap.}
The script \texttt{scripts/check\_attempt\_batch\_graphs\_2026\_10\_01.py} checked the $k=3$ construction on all 6,720 ordered five-vertex paths and all 1,680 ordered four-cycles of $K_8$. Each uses at least three colours. The script uses only the Python standard library; these finite checks supplement the symbolic proof.

The progression argument used equality in the integers. Replacing integers by residues can introduce nonzero elements killed by two, invalidating the four-cycle calculation; no modular version is asserted. The palette count includes unused sums, so it remains a valid upper bound for arbitrary $n$, not only powers of two. The smallest cases $n=2,3,4$ are covered without relying on the displayed five-vertex obstruction.

For this method to give the desired linear estimate, it would suffice to construct arbitrarily large progression-free integer sets $A$ with $|A+A|=O(|A|)$. This attempt does not construct such sets or claim they exist. More generally one may need a different edge-colouring rule; the failure of one additive rule is not a lower bound for all star colourings.

\paragraph{Final status.}
The additive lemma, explicit family, and bound (1) are proved, together with a concrete failure of consecutive labels. The requested uniform linear bound is \textbf{UNRESOLVED} by this attempt.

\subsection{Sources}
\begin{itemize}
\item[S1] ULAM.AI and the UnsolvedMath Contributors, supplied problems.json, record 3116 (OPG-37275), OpenGarden collection. Catalogue identity and source transcription; CC BY 4.0.
\url{https://huggingface.co/datasets/ulamai/UnsolvedMath}.
\item[S2] Open Problem Garden, Star chromatic index of complete graphs. Original problem and discussion; the mathematical formulation below is expanded from the supplied source record.
\url{http://www.openproblemgarden.org/op/star_chromatic_index_of_complete_graphs}.
\item[R1] Z. Dvořák, B. Mohar and R. Šámal, \emph{Star chromatic index}, arXiv:1011.3376; Journal of Graph Theory 72 (2013), 313--326. Checked primary abstract for the established near-linear context.
\url{https://arxiv.org/abs/1011.3376}.
\end{itemize}
\end{document}
